\SourceLeaf{399}%
\[
\begin{array}{ll}
x+n\tfrac12(x'+x''-2x),&y+n\tfrac12(y'+y''-2y)\\
x'+n\tfrac12(x''+x-2x'),&y'+n\tfrac12(y''+y-2y')\\
x''+n\tfrac12(x+x'-2x''),&y''+n\tfrac12(y+y'-2y'')
\end{array}
\]
oder, wenn man
\[
\begin{aligned}1-n&=\alpha\\\tfrac12n&=\beta\end{aligned}
\]
setzt,
\[
\begin{array}{ll}
\alpha x+\beta(x'+x''),&\alpha y+\beta(y'+y'')\\
\alpha x'+\beta(x''+x),&\alpha y'+\beta(y''+y)\\
\alpha x''+\beta(x+x'),&\alpha y''+\beta(y+y')
\end{array}
\]
\par
Von den Punkten \(C,C',C''\) fällt man Senkrechte auf \(A'A'',A''A,AA'\) und sucht die Lage der drei Schnittpunkte dieser Senkrechten. Die Koordinaten des Schnittpunkts der beiden letzten Senkrechten seien
\[
\xi,\quad\eta
\]
Sie lassen sich mithilfe des folgenden Lehrsatzes bestimmen.

Wenn
\[
\begin{array}{ll}
a,&b\\a',&b'\\a'',&b''\\a''',&b'''
\end{array}
\]
die Koordinaten von vier Punkten sind und die Gerade durch den ersten und zweiten Punkt senkrecht auf der Geraden durch den dritten und vierten Punkt steht, gilt
\begingroup\small
\[
\begin{aligned}
\tfrac{b'-b}{a'-a}&=\text{Tangente der Neigung der ersten Geraden gegen die Abszissenachse}\\
\tfrac{b'''-b''}{a'''-a''}&=\text{Tangente der Neigung der zweiten Geraden gegen die Abszissenachse}
\end{aligned}
\]
\endgroup
Da die eine Neigung um \(90^0\) größer als die andere ist, ist das Produkt der beiden Tangenten folglich \(=-1\), also
\[
\tfrac{b'-b}{a'-a}.\tfrac{b'''-b''}{a'''-a''}=-1
\]
In unserem Fall gilt also
